% page 410 du fac-simile (image p-409 du scan)
% bloc 152.4 x 233.3 mm ; marge gauche 8.0 mm
\pgc{-12.700mm}{0.000mm}{
\l{\cel{3}402.}
\l{}
\l{\sou{obturar}. (\sou{trans., per})\cel{2}(\sou{tekn.)}\cel{3}Stopar.}
\l{}
\l{\sou{obtuza.}\cel{1}I. Di qua la angulo esas rondatra, senpintigita.}
\l{\cel{3}- II. (\sou{geom.)}\cel{3}Angulo plu granda kam orto. - III.}
\l{\cel{3}(\sou{metaf.}) Spirito-nesagaca. - EFIS.}
\l{}
\l{\sou{obuso}. (\sou{tekniko milit}.) Projektilo kava, oblonga, sen anso}
\l{\cel{3}e sen kulo, qua spliteskas per perkuteso, lansate per}
\l{\cel{3}kanono. - FI.}
\l{}
\l{\sou{oceloto.}\cel{2}(\sou{zool.)}\cel{4}Kato-tigro di Amerika. - L.\cel{1}\sou{felis}}
\l{\cel{3}\sou{pardalia}. - DEFIRS.}
\l{}
\l{\sou{oceano.}\cel{2}Vasta extensajo de aquo saloza, qua okupas gran-}
\l{\cel{3}da parto di la Terglobo. - DEFIRS.}
\l{}
\l{\sou{ociar}. (\sou{netrans}.)\cel{2}Agar nulo, facar nulo, pro repugneso da}
\l{\cel{3}laboro. - FIS.}
\l{}
\l{\sou{ocidar}\cel{1}. (\sou{trans.})\cel{2}Mortigar vole e violente. - EFI.}
\l{}
\l{\sou{ocidento}. (\sou{literaturo}). Loko di la horizonto ube suno su}
\l{\cel{3}kushas. - II. Parto di la Terglobo, regiono situita}
\l{\cel{3}ne for la regiono ube suno semblas kushar su. - DEFIS.}
\l{}
\l{\sou{ocilar}. (\sou{netrans}.)Deviacar de sua baricentro, e retrovenar}
\l{\cel{3}ad ica per movado alternanta konstanta. -\cel{2}DEFIS.}
\l{}
\l{\sou{ocipito}. (\sou{anat.)}\cel{2}La infrajo di la dopa parto di la kapo.}
\l{\cel{3}- EFIS.}
\l{}
\l{\sou{ocitar.}\cel{2}(\sou{netrans}.)Apertar sua boko ye movo spama di expiro,}
\l{\cel{3}sequata da expiro plu longa, produktita da la bezono dor-}
\l{\cel{3}mar, da hungro, da fatigeso, da enoyo, da tedeso. - L.}
\l{}
\l{od.\cel{2}- Konjunciono qua expresas alternativo.}
\l{}
\l{odo. I. (en\cel{1}\sou{la epoki antiqua}). Poemo lirika, dividita aden}
\l{\cel{3}strofo, antistrofo ed epodo, quan la koro kantis evolu-}
\l{\cel{3}cionante. - II. Poemo lirika, dividita aden strofi inter-}
\l{\cel{3}simila pri nombro e la mezuro di la versi. - DEFIRS.}
\l{}
\l{\sou{odalisko}. I. Sklavino en haremo, destinita por servar la}
\l{\cel{3}spozini di la sultano. - II. Muliero ek haremo. - DEFIRS.}
\l{}
\l{\sou{\textquotedbl{}odeono\textquotedbl{}}.\cel{2}(en\cel{1}\sou{epoki antiqua)}.\cel{2}Edifico en qua onu repetis}
\l{\cel{3}la muzikaji destinita a kantesor en la teatro. - DEFIRS.}
\l{}
\l{\sou{odiar}.\cel{2}(\sou{trans.)}\cel{1}I. (\sou{ulu).}\cel{1}Malvolar profunde. - II. (\sou{ulo).}}
\l{\cel{3}Repugnesar profunde.da (lu). - EFIS.}
\l{}
\l{odiseo. Naraco di voyajo ed aventuri diversa. - DEFIRS.}
}
